% ============================================================
%  CHAPTER V
%  Reverse Perspective
%  Source pp. 93–117
% ============================================================

\chapter{Reverse Perspective}
\markboth{Chapter V. Reverse Perspective}{Chapter V}

\srcpage{93}

The term ``reverse perspective'' is sufficiently vague, and different authors give it
different meanings.  The usual meaning is an increase in the depicted dimensions of
the distant parts of objects compared with the near parts.  Sometimes the term
``reverse perspective'' is applied to the entire system of depicting pictorial space
on the picture plane in Old Russian painting, to the totality of conventions
characteristic of it, including those that have no direct relation to a perspective
system.  And indeed, Old Russian painting is characterised not only by departures
from the linear perspective familiar to the modern viewer, but also by other
distortions of the visible picture.  In the present book the term ``reverse
perspective'' is used in a purely geometrical sense, irrespective of the cause of
its appearance.

Reverse perspective in medieval art, and specifically in Byzantine and Old Russian
art, has been the subject of numerous studies.  First one should mention the works in
which reverse perspective is explained as the artist's incompetence or error.  Among
these one may cite the monographs of N.\,A.~Rynin, where he speaks of reverse
perspective as an ``erroneous technique''.\footnote{N.\,A.~Rynin.
  \emph{Nachertatel'naya geometriya.  Perspektiva.}  Petrograd, 1918, pp.~84--85;
  the same author, \emph{Materialy k istorii nachertatel'noy geometrii.}  Leningrad,
  1938, p.~85.}
Although this explanation is undoubtedly the simplest, it cannot be regarded as
scientific.

\srcpage{94}

Works that give a scientific explanation for reverse perspective can be divided
roughly into two classes.  To the first class belong those in which the foundations
of reverse perspective are derived from the specific character of medieval art, from
its symbolism, philosophical content, and tendency toward the ``timeless'' and
``transcendent,'' and so forth.  In particular, many scholars link the peculiarities
of Old Russian icons and frescos to religious dogma.  According to this view, the
medieval master depicted on icons another, unreal world, distinct from the earthly
one.  This is emphasised by the gold background conveying the sky (instead of the
natural blue), by the fantastically unreal depiction of objects (icon rocks), and so
on.  In order to show an ideal world not subject to earthly laws --- in particular
the laws of visual perception --- an artificial, reverse, perspective was introduced
in place of the natural one.  It must be said that such a point of view is not
without interest and has every right to exist.  Several formal features of medieval
painting are indeed connected with motives of this kind; it seems, however, that the
significance of philosophico-theological ideas has been greatly exaggerated here.
This is attested in particular by the presence of elements of reverse perspective in
antique art, in the art of the medieval East, in children's drawings, and so on.

To the same class of explanations one may assign the highly far-fetched supposition
that the artist painted the icon from a kind of interior vantage point, as if from
the position of our \emph{vis-à-vis}, imagined as being placed inside the picture.
This logical construction, going back to O.~Wulff, was thought fit to be supported
in one recent work by B.\,A.~Uspensky.\footnote{B.\,A.~Uspensky.  O semiotike
  ikony~[On the semiotics of the icon].---\emph{Trudy po znakovym sistemam}, V.
  Tartu, 1971, p.~199.}

To a different class of works belong those in which reverse perspective is connected
not with the artistic and substantive nature of medieval painting, but with certain
objective laws arising from the process of visual perception.  One should mention
first of all the work of A.\,V.~Bakushinsky,\footnote{A.\,V.~Bakushinsky.
  Lineynaya perspektiva v iskusstve i zritel'nom vospriyatii real'nogo prostranstva
  [Linear perspective in art and in visual perception of real space].---\emph{Iskusstvo},
  1923, No.\,1, pp.~213--261.}
his ingenious attempt to explain the appearance of reverse perspective by the
binocular nature of human vision.  In his view, each eye of a person sees the world
according to the laws of linear perspective, but the superposition of these two
images upon each other leads to the effect of reverse perspective.  The diagrams and
calculations he adduces show that this is indeed the case for objects whose
dimensions are smaller than the interpupillary distance and which are placed
sufficiently close to the observer.  These limitations make Bakushinsky's explanation
hardly effective, since such painting would depict only objects no bigger than a
matchbox, examined from a distance of 10--30\,cm.  But Bakushinsky does cite in his
work a series of interesting observations not directly connected with his theory
(which will be partly discussed below).

Another work of the same class is the monograph of L.\,F.~Zhegin.\footnote{L.\,F.~Zhegin.
  \emph{Yazyk zhivopisnogo proizvedeniya (uslovnost' drevnego iskusstva)}
  [The language of the pictorial work (the conventions of ancient art)].
  Moscow, 1970.}
He considers that with sufficient justification one may take human vision to be
monocular; at each moment the viewer sees the world by the strict rules of linear
perspective; but as he moves about and shifts the position of his viewpoint he sees
several different pictures and, mentally summing them up, creates on the picture
plane a certain synthetic image possessing properties of reverse perspective.  Thus,
if Bakushinsky reduced everything to binocularity, Zhegin regarded the mobility of
the viewpoint as the sole justification of reverse perspective.

Both authors just named, despite their ingenuity and a number of interesting
observations, make the same mistake.  They start from the naive assumption that a
person sees his retinal image (the untenability of this proposition was demonstrated
in Chapters~II and~III).

\srcpage{95}

As for the monograph of L.\,F.~Zhegin, its specific character requires some
additional comment.  Individual parts of the work are very unequal in quality.  When
the author shares his thoughts as a painter, sharing his experience of the living
perception of works of art, his observations are interesting.  But where he appears
as a geometer his constructions are below all criticism, and the results he obtains
are erroneous.  The muddled geometrical arguments of Zhegin are naturally perceived
as incomprehensible, but representatives of the humanities assume that the
incomprehensibility is due to the mathematical character of the material presented.
Those who think so should be disappointed --- these arguments contain no rational
geometrical content.  In essence, the only rational geometrical proposition of the
work under discussion is the almost banal assertion that if a person examines an
object from different sides and then attempts to show all of this simultaneously,
the effect of reverse perspective will arise.\footnote{%
  This can be illustrated by the example of the icon-rock (\emph{ikonno-gornoye})
  construction, which Zhegin himself and supporters of his views regard as a sort of
  quintessence of the theory (see pp.~29, 90 and fig.~33 of Zhegin's monograph, and
  also pp.~194, 216 of the above-cited work of B.\,A.~Uspensky).  It is claimed that
  ``icon rocks'' depict the horizon, or more precisely its perspective distortions in
  the system of ``active'' space characterised by a dynamic viewpoint.  But the
  dynamics of the viewpoint (the artist's movement) relative to the horizon can
  contribute nothing, since the horizon in perspective theory is an infinitely remote
  region of space, and any finite displacements of the artist, as mathematics teaches,
  are equivalent to his standing still.  The very formulation of the problem therefore
  turns out to be mathematically meaningless.

  It seems absurd too from the standpoint of plain common sense.  It is hard to
  imagine a greater dynamics of viewpoint than sailing from Europe to America.  If
  something were to change in the region of the horizon, people would long since have
  described such miracles.  The sea horizon, however, is universally associated with
  prolonged and tedious monotony.

  Although the formulation of the problem seems meaningless, let us follow the
  author's logical constructions further.  As the next step he introduces ``curved
  lines of sight,'' without explaining what these are or why these lines should be
  curved.  The completely arbitrary introduction of new concepts without explanation
  or definition is very characteristic of Zhegin.  This may be admissible when an
  author wishes to enrich the imagery of his language, but it is absolutely
  impermissible when such concepts are used to obtain geometrical conclusions.  If the
  author of the icon-rock construction is unable to explain the reason for the
  curvature of the lines of sight, he ought at least to have explained where they are
  curved~--- in objective space or in the artist's consciousness.  Since the author of
  the theory dismisses these ``trifles'' with astonishing ease, it becomes necessary
  to examine both cases.  If the lines of sight are curved in real space, in the
  atmosphere, then the whole chain of reasoning is valid and a displacement of the
  image of the observed point in visual perception relative to its true position would
  indeed arise, since human perception is ``accustomed'' to light travelling in
  straight lines; the perceived position of the point would lie on the line going from
  the eye in the direction from which the ray of light ``entered'' the eye, since the
  brain cannot sense the curvature of the propagation of light rays.  This phenomenon
  is called refraction and is well studied.  However, the refraction needed for
  Zhegin's constructions has never been observed, and any physicist can readily
  explain that in the earth's atmosphere it cannot exist.  If one supposes that the
  lines of sight are ``curved'' not in the atmosphere but in human consciousness, then
  no ``rotated'' direction of ``entry'' of a ray into the eye exists, and the
  refraction picture loses its meaning.  Thus, whichever of the two possible
  assumptions one adopts, the reasoning leading to the stated geometrical constructions
  turns out in both cases to be erroneous.

  \smallskip
  Let us nevertheless agree provisionally to disregard this circumstance.  Following
  further the logic of the proposed theory of the icon-rock construction, one sees
  from fig.~33 of Zhegin's monograph that an initially continuous line would
  ``break'' into separate segments.  But the reasoning that led to these breaks holds
  not only for the points marked in the diagram but for all points between them.
  Then, as mathematics teaches, what will occur is not ``breaks'' but only
  ``stretchings.''  In that case the further consequences of the erroneously obtained
  ``breaks,'' and with them the whole theory, lose their meaning.

  As may be concluded from the foregoing, each of the first three fundamental steps
  in the theory of the icon-rock construction is untenable; even one of these three
  errors would be sufficient to invalidate everything that follows~--- especially
  since the subsequent reasoning also contains mathematical inconsistencies.  The
  other geometrical constructions of Zhegin, not directly connected with the
  icon-rock construction, are for the most part no stricter.  But then his
  conclusions about shifts, fractures, curvatures, and constrictions of the field of
  view, and so on, have no rational geometrical meaning.  At best they may be
  interpreted as a painter's figurative description of his perception of icons,
  making no claim to explain all these effects.}

\srcpage{96}

Somewhat apart stands the extremely interesting work of P.\,A.~Florensky, written in
1919 and published posthumously in 1967.\footnote{See:
  \emph{Trudy po znakovym sistemam}, III.  Tartu, 1967.}
It represents the first part of an intended study on reverse perspective.
Unfortunately the author did not realise his plan in full.  The published part
contains a rich historical survey, a critical examination of the widespread belief
that linear perspective is the only admissible system, and the suggestion that
depending on artistic goals the perspective system may vary.

The cited article of P.\,A.~Florensky criticises the theoretical starting positions
that confine our ideas about spatial constructions to linear perspective and
underlines their artificiality.  The analysis of reverse perspective and other
questions connected with the system of constructing the image on the picture plane in
Byzantine and Old Russian art were presumably to be presented in the second part of
the work, but it was never written.  Therefore P.\,A.~Florensky's point of view on
the reasons for the appearance of reverse perspective in Old Russian art remained
unknown.\footnote{Florensky himself writes as follows: ``The author has no intention
  at all of building a theory of reverse perspective'' (op.\,cit., p.~414).}

E.~Panofsky in his well-known work ``Perspective as Symbolic
Form''\footnote{E.~Panofsky.  Op.\,cit.}
holds that in the process of the historical development of art different perspective
systems appear with changes in ideology.  According to Panofsky (who says not a word
about reverse perspective as such), antiquity did not know linear perspective,
although it intuitively approximated it, and was unable to convey unified space,
confining itself to the volumetric depiction of individual objects.  Medieval art,
proceeding from new ideological tasks, began to depict unified space, though in doing
so it reached a complete ``destruction'' of perspective.

\srcpage{97}

Only the masters of the Renaissance found methods for perspectival depiction of space
and the objects disposed in it as a unified whole.  Panofsky connects this process
with changes in philosophico-theological views in the transition from antiquity to the
Middle Ages and from the latter to the Renaissance.

\medskip
\noindent
$* \quad * \quad *$
\medskip

Let us attempt to examine the problem of the appearance in medieval art of images
characterised by reverse perspective, relying on the peculiarities of human spatial
perception (see the preceding chapters).

First let us agree that purely geometrical properties of the image will be
characterised.  By \emph{direct perspective} we shall mean a manner of depicting
extended objects in which their characteristic linear dimensions progressively
decrease as distance from the viewer increases; by \emph{reverse perspective} --- the
opposite manner of depiction, in which characteristic linear dimensions increase with
increasing distance.  Clearly the axonometric system of depiction constitutes an
intermediate case, on the boundary between direct and reverse perspective.  This
suggests that even an arbitrarily ``weak'' influence can transform the axonometry of
the close foreground into reverse perspective.  On the other hand, Old Russian and
Byzantine painting had as its chief subject shallow, near spaces, and therefore the
analysis of possible transformations of axonometric images of near objects is the
very basis capable of explaining the appearance of reverse perspective in Old Russian
art.

Analysis of Old Russian and Byzantine painting shows that the appearance of reverse
perspective is not connected with the predominant action of any single factor.  It is
remarkable that a whole series of apparently quite heterogeneous causes leads in the
end to the appearance of elements of reverse perspective in the image.  Therefore the
characteristic attempts of previous works to find a single explanation for reverse
perspective in medieval art inevitably impoverished the investigation.

\srcpage{98}

The various causes to be examined below transformed the axonometrically based image
of the near foreground both in the direction of direct and of reverse perspective.  In
the process, owing to the objective circumstances discussed below, the number of
transformations in the direction of reverse perspective in medieval art proved to be
significantly larger than the number in the direction of direct perspective.  This is
what led to reverse perspective becoming one of the most striking and conspicuous
features of this painting.

Transformations of the axonometrically based image occur, as already noted, for
different reasons; they conveniently fall into two categories --- those directly
connected with the process of human visual perception of space, and those not
connected with this process.  To the second category we assign, for example, the
artist's use of reverse perspective for compositional reasons (such questions will be
the subject of the next chapter).

Examining in the present chapter only what is directly connected with visual
perception, let us clarify certain aspects of the processes by which the retinal
image is processed by the visual perception system.  These processes, let us recall,
concern chiefly the comparatively narrow and more important field of sharp vision.
Consequently the results obtained below will be valid only under these conditions.
This limits the scope of the investigation to some extent (in subsequent chapters the
question will be examined in a broader formulation).

Returning again to the mechanisms of the visual perception process and their
reflection in perspective systems, we shall additionally take into account those
aspects of this process that were omitted in the simplified treatment of Chapter~III.
It is precisely these refinements that explain the ``deformations'' of axonometric
images of the close foreground that proved to be the most characteristic feature of
the first-approximation perceptual perspective system.

We shall adopt the following sequence of refinements: 1~--- allowance for the
phenomenon of ``superconstancy''; 2~--- perception of objects in foreshortening;
3~--- allowance for the non-Euclidean character of perceptual space; 4~--- allowance
for the form-constancy mechanism.

\medskip
\noindent\textbf{1.}\quad
``Superconstancy'' sometimes arises under binocular vision.  The monocularity assumed
above is a reasonable idealization of the vision process, since the distance between
the eyes is, as a rule, small in comparison with the depicted objects and the
distances to them.  Therefore allowance for the actual binocularity of human vision
would seem not to affect substantially the depiction of an object on the picture
plane.

\srcpage{99}

As modern experiments show (see Appendix~5), binocularity affects the quantitative
side of the size-constancy mechanism.  All authors who have conducted the relevant
experiments unanimously note that binocular vision is characterised, compared with
monocular (all else being equal), by a more strongly expressed constancy, i.e.\ in
the result of the action of this constancy mechanism the relative enlargement of the
sizes of distant objects turns out to be stronger under binocular vision than under
monocular vision.  In some cases the effect of ``superconstancy'' was registered,
when under binocular vision a distant object appeared larger than its true size.

If one proceeds from this experimental fact, the following scheme of binocular
vision may be proposed: suppose that under monocular vision the near foreground is
axonometric; then under binocular vision distant objects (within this foreground)
will appear larger than their true size, whereas objects in the immediate vicinity of
the observer are always seen in true sizes.  This will lead to the appearance of an
effect of reverse perspective.  Continuing this thought experiment, as the distance
to the observed objects increases still further they will begin to decrease, since
in the end any object at an infinitely large distance will be seen as a point.
Straight path in this case will have the appearance shown in Ill.\,36.  Comparing
the image given with the right-hand part of Ill.\,17, one confirms that the type of
perceptual image has changed only within the near foreground, where axonometry has
given way to reverse perspective.

\figblock{36}{Perceptual perspective of a path taking into account the natural
  transformations of the foreground.}

Numerical estimates (based on the experiments mentioned above and given in
Appendix~5) give grounds to expect such an effect for not too wide objects situated
at distances of 3--6\,m from the observer.  The magnitude of the reverse perspective
may be of the order of 2°, i.e.\ it should be classified as very weak reverse
perspective.

\medskip
\noindent\textbf{2.}\quad
\emph{Perception of objects in foreshortening.}  The appearance of reverse
perspective as a result of ``superconstancy'' gives a very simple explanation for
this type of geometry of visual perception of near regions of space.  Unfortunately,
experiments in the psychology of visual perception make clear that ``superconstancy''
is not too common.  One may suppose that in the past there were more people
possessing this perceptual peculiarity (which is not excluded), but such a hypothesis
is difficult to confirm convincingly.

However, there exists yet another possibility for the appearance of reverse
perspective as a result of the action of the size-constancy mechanism~--- one not
based on ``superconstancy'' and therefore capable of manifesting itself in practically
every person (see Appendix~6).  As theoretical analysis shows, the transition from
axonometric perception or even from visual perception possessing properties of slight
direct perspective, to vision in slight reverse perspective, is possible when the
foreshortening of an object changes.  To illustrate the nature of such a
transformation, let us turn to Ill.\,37.  The left column of the figure

\srcpage{100}

schematically shows the process of transition from linear perspective (the retinal
image) to perceptual perspective (the geometry of perceptual space) for the already
examined case of a narrow path.  To simplify the constructions, the assumption is
made that only the transition from the retinal image to the geometry of perceptual
space for the near region of space is shown, and for a person whose vision is
characterised by full constancy.  Diagram~A gives the corresponding retinal image.
Under the action of the size-constancy mechanism the retinal image will be stretched
in the directions shown by the arrows, producing diagram~B, where the edges of the
path are already seen as parallel.  However (as follows, for example, from Ill.\,17)
the stretchings occur not only horizontally but also vertically.  This process too is
explained by arrows, and leads to diagram~C.  Diagrams~B and~C differ only in that
in the second the shown section of the path has become more extended vertically.  A
quite different result can be produced by the same transformations if the path is
directed not away from the observer but ``diagonally.''  Diagram~D gives the retinal
image for the same path with an altered foreshortening.  The size-constancy
mechanism, by stretching in the horizontal direction, will, as in diagram~B, make its
edges parallel, leading to diagram~E.  As for the vertical stretchings that in
diagram~C merely elongated the image of the path without changing its character,
they will now lead to diagram~F, where the path is seen in clear reverse perspective.
This will happen because the far edge of the path will rise (stretch) more than the
near edge; this is schematically shown by arrows of different lengths in diagram~E.
The same phenomenon is visible in diagrams~D and~B, where the distant parts also
stretch more than the near parts.  It is precisely this non-uniformity of stretching
that will lead to the appearance of reverse perspective.  This effect could not
manifest itself in the diagrams of the left column of Ill.\,37, since the edges of
the path received no lateral displacements (the lines representing the edges did not
rotate), so the non-uniformity of vertical stretching, which also exists in this
case, could not change the geometric character of the appearance of the path.

\figblock{37}{Transformation of the visible appearance of a path observed in
  foreshortening.}

\srcpage{101}

As is clear from this brief and highly schematic description, to see in slight
reverse perspective requires no vision with certain peculiarities such as
``superconstancy,'' nor binocularity, since full constancy is not infrequently
observed even under monocular vision.  Moreover, reverse perspective in diagram~F can
appear even if the lines in diagram~E are not strictly parallel but slightly
converging with depth, so long as the difference in the stretching of the nearest and
more distant regions (arrows in diagram~E) is sufficiently pronounced.  It follows
that nearly everyone can see in slight reverse perspective today.  Returning to the
question of binocularity, one may assert that it enhances the necessary effects here
as well, since, as already noted, binocular vision is characterised by more strongly
expressed constancy.

If the reader wishes to convince himself of this, the following simple method may be
suggested.  One should stand in a room with a clearly pronounced linear structure of
the floor (for instance, parquet strips), facing a blank wall, and clear the space in
front from all objects.  One should stand approximately 4--5\,m from the wall.
Facing the wall and looking straight ahead, one should find a spot where the strips
going away from the viewer appear exactly parallel (this typically corresponds to a
distance of the order of 3--5\,m).  After this one should shift one's gaze to the
side, by approximately 20--50°, to the same distance, and the linear structure of
the floor will then be directed diagonally relative to the line of gaze.  Under
these conditions it will be seen, within the region of sharp vision, in slight
reverse perspective.  It is important only that the width of the strips at which the
gaze is directed be neither too large nor too small; depending on the distance to the
viewed region one may regard a set of several strips together.  It should perhaps be
added that effects of this kind are best observed in small closed spaces.

If someone is unable to see the strips in slight reverse perspective under the
indicated conditions, this will only mean that he is very well and firmly ``trained''
by habitual linear perspective.

The calculations given in Appendix~6 indicate that under the conditions described
here, natural reverse perspective in visual perception can be of the order of 10°,
i.e.\ it is significantly stronger than that associated with the phenomenon of
``superconstancy,'' while still remaining ``weak'' reverse perspective.

\medskip
\noindent\textbf{3.}\quad
\emph{The non-Euclidean character of the geometry of perceptual space} was discovered
in the mid-twentieth century.  Processing the numerical characteristics of human
visual perception, Luneburg interpreted its peculiarities in terms of non-Euclidean
geometry.  It turned out that in contrast to objective space, which is described by
ordinary school Euclidean geometry, perceptual space (more precisely, the region of
it close to the observer, where the regularities of binocular vision are decisive)
possesses properties of Lobachevsky space.  (For a more detailed account and
literature references see Appendix~7.)

If one takes this experimentally confirmed result as the starting point of the
argument, one can in particular pose and solve the following geometrical problem (see
Appendix~7).  Suppose a rectangle is drawn on some plane, for example on the floor.
What properties will the figure drawn on the floor have in a person's visual
perception?  Bringing the Euclidean floor plane and the Lobachevsky plane that has
arisen as the image of the floor in the person's consciousness into proper
correspondence, one can verify that the depicted rectangle acquires in visual
perception the form of a so-called Lambert quadrilateral.

\srcpage{102}

Without describing all the properties of such a quadrilateral, let us note only that
its side more distant from the observer will be greater than the nearer side.  But
this \emph{is} ``reverse perspective,'' for in reality the near and far sides of the
rectangle are equal, while in visual perception it turns out to be expanding in
depth.

This conclusion shows that real vision in reverse perspective does not at all require
that the person's vision depart from the norm, as one might have concluded from the
first point, where the effect of vision in reverse perspective was connected with the
comparatively rarely observed phenomenon of ``superconstancy.''

Thus, the phenomenon of ``superconstancy,'' the non-Euclidean character of perceptual
space geometry, and especially the visual transformations associated with viewing
objects in foreshortening, explain reverse perspective in art by entirely natural
causes.  One should add only one supplement to what has been said: visual effects of
this kind are especially well observed in closed spaces.  This is not the place to
interpret this fact from the standpoint of the mechanisms of visual perception; it
suffices simply to note this experimentally confirmed circumstance.  The importance
of such a supplement stems from the fact that it is precisely from reproductions of
``furniture objects'' (footstools, etc.) that one usually judges the characteristic
features of perspective.  Observing such objects in slight reverse perspective or in
axonometry, the artist working from memory often did not introduce any geometrical
distortions of natural perception.  If this is so, images connected with the
described vision must also exist in medieval art.  But the numerous objects that are,
for example, on icons depicted according to the rules of reverse perspective cannot
serve as clear proof of the correctness of the arguments developed above, since any
such case can be interpreted as a form of symbolism, or as a desire to subordinate
the image to the plane of the board, or in some other analogous way.  Nevertheless
much convinces us of the reality of the vision and depiction of the near foreground
according to the laws of reverse perspective.

\srcpage{103}

Above all, images of this kind appear universally at that stage of development of the
visual arts when artists attempt to convey the visible picture directly (unaware of
the existence of the intermediate image in the process of visual perception --- the
retinal image --- to which the system of linear perspective corresponds).

Indeed, not only Byzantine, Old Russian, medieval Armenian and Georgian art, and the
art of other countries connected with Byzantine culture, possesses this peculiarity.
Reverse perspective can be found in the medieval art of Iran, India, China, Japan,
Korea, and a number of other countries of entirely different culture.

This manner of depiction is also a characteristic feature of children's drawings
and, as special studies have shown, cannot be reduced to ``childish
ineptitude.''\footnote{A.~Amdgren.  \emph{Die umgekehrte Perspektive und die
  Fluchtachsenperspektive.}  Uppsala, 1971.  More profound research into the
  reverse-perspective features of children's drawings is contained in a number of
  works by I.\,P.~Glinskaya, whose final account can be found in her doctoral
  dissertation ``Formirovaniye sposobov ovladeniya prostranstvennoy informatsiey na
  ploskosti u mladshikh shkol'nikov'' (Leningrad, 1973).}

Such a universal and seemingly spontaneous emergence of reverse perspective must have
a common cause, and it is natural to see that cause in the regularities of human
visual perception.

To illustrate the foregoing, let us consider two examples.  Let us examine a
miniature from the Vatican Joshua Scroll (7th century), where the antique tradition
is felt quite distinctly.

\srcpage{104}

In this miniature the warriors stand in a picturesque group in varied and lively
poses, one of them standing even with his back to the viewer; the main personage is
not depicted frontally --- he has turned his head to address those kneeling before
him; all the figures are proportionate.  The scene contains no geometrical
distortions of the figures, despite the varied foreshortening.  All the more
remarkable, then, that the footstool of the throne is given in slight reverse
perspective (of a magnitude of the order of 5--6°).  Regardless of whether one
considers this depiction to be medieval or, with better justification, Hellenistic,
it is hard to imagine that in this illustration the artist decided, by distorting the
visible form of the footstool, to lend the whole image some symbolic meaning.  One
can say with complete confidence that the artist showed equally the figures of the
people and the footstool in the main as he was accustomed to seeing them in
nature.\footnote{V.\,N.~Lazarev, characterising this image, says that the crowd is
  depicted with impressionistic lightness.  He dates the prototype of the image to
  the Hellenistic style of the 5th century.  See: V.\,N.~Lazarev.
  \emph{Istoriya vizantiyskoy zhivopisi} [A History of Byzantine Painting], vol.\,I.
  Moscow, 1947, p.~55.}

\figblock{38}{Miniature from the Vatican Joshua Scroll.  7th century.}

As another example let us take the self-portrait of the Korean realist painter (18th
century) Kim Hong-do.  The low table shown at left and the objects on it are given in
slight reverse perspective.  The degree of reverse perspective~--- a divergence of
the lines of the order of 3--4°~--- is close to the analogous quantity in the
previous example.  The general character of this work also speaks for the conclusion
that the author was conveying space here as he saw it.

\figblock{39}{Kim Hong-do.  Self-Portrait.  Painting on paper.  18th century.}

\srcpage{105}

These two works, separated in both time and space, belonging to different epochs and
different cultures, speak to the existence of some common cause for the rendering of
extended objects in slight reverse perspective.  The only possible common cause is
natural vision in reverse perspective, whose features were discussed above.

It is worth emphasising once more that reverse perspective connected with the
artist's direct reproduction of the geometry of his visual perception is
distinguished by two features allowing one to separate natural reverse perspective
from reverse perspective arising from other, artistically justified, impulses.
First, natural reverse perspective can appear only in the depiction of near objects,
especially those shown in foreshortening, and second, its magnitude must not exceed
approximately 10° (as in the two illustrations given).  This allows one to refine the
meaning of the words ``I see it this way,'' which contemporary artists~--- not
infrequently turning again to the possibilities of reverse perspective~--- like to
use.

\medskip
\noindent\textbf{4.}\quad
\emph{Allowance for the action of the form-constancy mechanism.}  If the three causes
of natural reverse perspective given above are connected with transformations of the
retinal image, with its transformation as a result of the action of the
size-constancy mechanism, a different but equally natural cause consists of effects
connected with the form-constancy mechanism.

As already noted in Chapter~II, when contemplating a familiar object whose form is
known from prior experience, a person sees it as more nearly conforming to its true
form than is given in the retinal image of the same object.  This observation by
psychologists was confirmed by a series of experiments permitting quantitative
assessment of effects of this kind.  In experiments conducted by A.\,A.~Smirnov, the
visible size of a square plate was fixed as it was rotated at different angles to the
line of sight.\footnote{A.\,A.~Smirnov.  Zavisimost' konstantnosti vospriyatiya
  velichiny ob''yektov ot ugla povorota ikh k linii vzora nablyudatelya pri
  razlichnykh distantsiyakh nablyudeniya [Dependence of the constancy of size
  perception of objects on the angle of rotation relative to the observer's line of
  sight at various viewing distances].---\emph{Zritel'nyye oshchushsheniya i
  vospriyatiya.}  Moscow\,--\,Leningrad, 1935, p.~259.}
The plate was rotated about an axis parallel to one of the sides of the square, so
that one side always retained its size for any rotation of the plate.  The other side
of the square appeared smaller the more the plane of the plate deviated from the
direction perpendicular to the line of sight.  The experiment showed that in an
inclined position the plate always appears wider than its projection built by the
rules of linear perspective (i.e.\ relatively wider than its retinal image).  In
particular, up to rotation angles of the plate of the order of 40° it continued to
appear almost square, although in the retinal image the ratio of the sides was
already of the order of 0.7.

\srcpage{106}

Let us give a diagram illustrating the experiment described (Ill.\,40, above).  The
thin lines show the actual position of the plate relative to the line of sight (arrow
in the diagram), and the thick lines show the position the plate would have to take
for its retinal image to correspond to what is seen.  In other words, the thick lines
correspond not to the actual but to the visible angle of rotation of the plate.  If
the artist faces the task of drawing the visible size of the plate, he will show the
width of the plate varying with rotation not in accordance with the thin lines but in
accordance with the thick ones.  Consequently he will depict the plane of the square
as shifted in the direction of approaching the position perpendicular to the line of
sight, i.e.\ to the picture plane.  This partly explains the tendency of an artist
working in the system of perceptual perspective to spread the images of the surfaces
of volumetric objects flat on the picture plane.\footnote{The form-constancy
  mechanism sometimes leads to an interesting artistic effect in portrait painting, as
  a result of which the depicted person not only ``follows the viewer with his eyes''
  as the viewer moves but also ``turns'' his whole body relative to the frame of the
  picture.  Since this question has no direct relation to medieval art, it has been
  placed in Appendix~11.  The mathematical side of this problem is completely
  elementary.}

Problems connected with the action of the form-constancy mechanism are illustrated by
the example of depicting a cubic stool (Ill.\,40, below).  Depiction~A corresponds to
ordinary axonometry; its deficiency is the non-correspondence of the form of the
seating surface to the configuration seen by the artist.  As a result of the
constancy mechanism, the person sees this form as more nearly conforming to the real
(square) form, and this is shown in depiction~B.  The visible configuration
corresponds to some fictitious rotation of the seating surface (dashed line).  If the
seating surface is depicted as shown in diagram~B (i.e.\ as it is actually seen), the
question arises of how correctly to show the legs of the stool.  The point is that
the form-constancy mechanism acts in relation to material objects having a definite
form, known to the viewer in advance.  This applies fully to the seating surface,
while the mental square at whose corners the legs touch the floor is not something
material and tangible, but a geometrical abstraction to which the form-constancy
mechanism of objects does not extend.  But then the contact points of the legs with
the floor will not shift, and depiction~B arises~--- giving an example of reverse
perspective, for the front legs are shorter than the back ones.\footnote{The trivial
  case of observing the stool from above, at a very small distance, when the back
  legs would be larger than the front ones even in a photograph, is not considered
  here.}

\figblock{40}{Diagrams for depicting a cubic stool.}

\srcpage{107}

To show that a person really does see the seating surface of the stool according to
diagram~B, an experiment was conducted in which a subject observed a stool with a
white seating surface and was required to select, from many parallelogram images
presented to him, the one corresponding to the visually perceived form of the seating
surface.  The selected parallelogram image was then placed next to the stool and
photographed.

\figblock{41}{Results of the experiment on determining the visible form of the seating
  surface of a stool.}

The appearance of the reverse-perspective effect as a result of the form-constancy
mechanism is not something specifically inherent to this mechanism.  If the artist
were faced with the task of depicting the stool from the front and below~--- from the
viewpoint of a cat, so to speak~--- the form-constancy mechanism would have
transformed the axonometric image into a depiction possessing properties of direct
perspective.  Indeed, in this case too the seating surface (observed from below)
would have ``rotated'' in the direction of approaching the plane perpendicular to the
line of sight, and as a result the back legs of the stool would not have lengthened
but shortened.  Thus the form-constancy mechanism leads to the effect of reverse
perspective for a certain direction of the line of sight in the case where the
picture is painted when viewed from the top and front.  In Old Russian painting
virtually all objects (footstools, seats, tables, etc.) are depicted when seen from
the top and front, and it is precisely this that has a decisive effect on the
appearance of the effect described above.

One may imagine the case when in depicting the stool it is important to show the legs
as equal in length; in that case depiction~C arises (Ill.\,40), in which the back
legs have ``lifted off the floor'' (thick lines).  Various intermediate modes of
depicting the stool are of course also possible.

The example of depicting the stool is interesting for illustrating the impossibility
of adequately depicting the visible spatial image on a plane (see Chapter~III).
Indeed, all three depictions in Ill.\,40 are unsatisfactory from this point of view:
in the first the configuration of the seating surface is incorrectly shown, in the
second the legs have acquired different lengths, in the third the legs have lifted off
the floor.  The choice among these schemes is left to the artist, and he makes it on
the basis of his artistic intentions.  Most commonly encountered is scheme~B.  It was
used already in antique art and occurs in practically the whole of medieval painting,
in particular Byzantine and Old Russian.  We shall confine ourselves here to a
reference to the icon ``The Nativity of St.~John the Forerunner'' (Ill.\,27), where
the legs of the table, stool, and cradle possess the property described.  Schemes~A
and~C occur as well, though much more rarely.

\srcpage{108}

Digressing somewhat from the question under discussion, let us draw attention to the
fact that medieval art showed a clearly determinate preference for scheme~B, whereas
a contemporary person would without doubt choose the axonometric scheme~A.  Evidently
the medieval artist could not imagine that upon the narrow strip of the seating
surface~--- corresponding to scheme~A~--- a saint could be accommodated with comfort
and dignity; at the same time the reverse-perspective effect arising from scheme~B,
for the reasons set out above, aroused in him no sense of protest, which would almost
inevitably appear in a modern viewer.  Although the modern person too would see in
nature the seating surface of the stool (owing to the form-constancy mechanism) as
more nearly square than in scheme~A (i.e.\ as in scheme~B), he has from childhood
been accustomed to perceiving pictures in the system of linear perspective.  The
method of compensating for the distortions characteristic of linear perspective
during the viewer's unconscious processing of the picture was discussed in detail in
Chapter~III.  Therefore the modern viewer ``feels'' that the strip of the seating
surface in the picture is larger than it appears, whereas the medieval viewer, who
knew only the system of perceptual perspective (which has the property of depicting
objects in such a way that subsequent unconscious transformation is not only
superfluous but harmful), would simply ``see'' it as it is drawn.  Moreover, the
modern viewer is accustomed to direct perspective and finds reverse perspective
unnatural.

This suggests that medieval painting was perceived in its own time somewhat
differently from the way modern people perceive it, even in a formally visual sense,
and this is worth bearing in mind.

Even today the unnaturalness for the near foreground of images built according to the
laws of linear perspective constantly leads artists to deviate from it and to turn to
modes of conveying spatiality more in keeping with the regularities of visual
perception.  Among the many violations of the system of linear perspective one can
find images constructed according to the medieval scheme~B.  Interesting in this
connection is a canvas by a pupil of K.~Briullov, one who undoubtedly passed through
a rigorous school of perspective training~--- the picture by A.\,S.~Mikhailov, ``The
Secretary Playing the Guitar.''  The vanishing point of objectively parallel lines on
the table reproduced in the said canvas~--- the left edge of the table and the line
through the points where the left legs touch the floor~--- lies closer to the viewer
than the table, exactly as with the analogous ``vanishing point'' in Ill.\,40B.  The
artist skillfully conceals this circumstance by placing the back leg of the table in
deep shadow.  He apparently understands that he is acting ``contrary to convention,''
but the medieval scheme allows him to solve the artistic problem before him, which
naturally had nothing to do with a deliberate geometrical distortion of vision for
formalistic reasons.  Consequently the medieval modes of depiction are not so strange
as they might appear at first glance.

\figblock{42}{A.\,S.~Mikhailov.  \emph{The Secretary Playing the Guitar.}  1851.
  State Russian Museum.}

\srcpage{109}

In general, when depicting near regions of space, elements of reverse perspective
appear on artists' canvases more often than one might think.  If one takes still
lifes, especially those painted after the mid-19th century, the far part of the table
on which the objects are arranged is often raised.  If the artist tried to depict the
legs of such a table, he would have been forced to make the back ones larger than the
front ones.  Such phenomena are entirely natural; they speak of the striving to
convey the visible image undistorted.

This natural pull of artists toward reverse perspective was expressed by K.\,F.~Yuon
in the following words: ``Anyone who does not know the laws of perspective theory
will almost necessarily depict objects in the reverse manner, as was done
systematically in all cases of Eastern and ancient folk art.''\footnote{%
  \emph{K.\,F.~Yuon ob iskusstve} [K.\,F.~Yuon on Art], vol.\,I.  Moscow, 1959,
  p.~47.}
To these words of Yuon one should add that this speaks not so much of ``ignorance of
perspective'' as of an intuitive striving for the direct transmission of visible
geometry.

\srcpage{110}

When the retinal image is transformed under the influence of the form-constancy
mechanism in the direction toward the true form, it is not only linear magnitudes
that undergo change but also angles.  This is visible, for example, from the
comparison of the forms of the seating surface (Ill.\,40).  The striving to depict
the angles in the image as they are seen can also lead to reverse perspective.

In this connection let us examine the depictions of a Gospel book and other
rectangular objects (Ill.\,43).  The most widespread mode of depicting a Gospel book
is one in which its two visible edges (fore-edges) are made of equal width; this
expresses the artist's wish to convey their actual equality.  Confining ourselves to
this case, the axonometric depiction of the Gospel book (Ill.\,43A) will be
characterised by three groups of parallel lines: vertical, horizontal, and inclined
to them at 45°.  The last is a consequence of the equality of the depicted edges,
which leads to the boundary between the two shown edges always having this direction.
As for the other visible boundaries of the edges, by axonometricity they must have
this direction as well.

The form-constancy mechanism, acting when a person contemplates real parallelepipeds,
leads to the visible magnitudes of the angles being significantly closer to right
angles than their retinal images.  In this case diagram~A has the defect that the
angles forming the silhouette of the upper edge of the Gospel book are too acute and
too obtuse (45° and 135°) compared with their visible magnitude.  In consequence of
this the medieval artist might decide to depict the outer upper-left angles $\alpha$,
$\beta$ and the lower angles $\alpha'$, $\beta'$ formed by the boards of the binding
and the edge boundaries, as more nearly right, leading to diagram~Ill.\,43B.  As for
the inner upper-right angles $\gamma$ and $\delta$, following the visible magnitude
of the angles would there lead to a break in the image, since simultaneously reducing
the angles $\gamma$ and $\delta$ (bringing them toward right angles) without creating
a break is impossible.  Medieval artists always considered the continuity of the image
more important than the correctness of lengths or angles, and therefore in depictions
of type~B always retained $\gamma = \delta = 135°$.\footnote{The artist watches
  angles $\alpha$, $\beta$, $\alpha'$, $\beta'$ and not $\gamma$, $\delta$, since the
  former determine the silhouette while the latter are unimportant in this sense.}
As a result in the depiction of the Gospel book the back board of the binding becomes
wider and taller than the front one, i.e.\ the effect of reverse perspective
arises.\footnote{The striving of the artist to preserve the continuity of the image
  (even at the cost of distorting the visible form) is entirely natural, since the
  human visual perception system is characterised first of all by the continuous
  character of the geometrical transformations in forming perceptual space from the
  retinal image.  These transformations (of the stretching and compression type) never
  lead to breaks in the visible space and objects disposed in it.  Naturally, never
  observing breaks, the artist considered it impossible to depict them, i.e.\ he
  proceeded from a kind of ``primacy of non-breakage.''  If on individual icons (very
  rarely) breaks in the image do occur, even a superficial analysis shows that these
  breaks have nothing to do with the system of visual space perception.

  It goes without saying that no mention is made here of cases where the artist for
  compositional or other reasons considers it necessary to show a certain object
  twice, or deliberately to divide it into parts and show each separately, etc.  What
  is meant is that each such individual object in the depiction (regardless of its
  properties in objective space) is shown as an unbroken whole, without structural
  absurdities.

  In the example of the stool depiction (Ill.\,40) this assertion can be illustrated
  thus: the artist will never show the legs of the stool according to scheme~A while
  showing the seating surface according to scheme~B, creating a break between the back
  legs and the seating surface, even though this would be a precise record of the
  numerical characteristics of visual perception.  One can show that the ``primacy of
  continuity'' has a strict mathematical justification and reduces to the requirement
  that the perspective system possess the property of commutativity (see
  Appendix~8).}

\srcpage{111}

Sometimes artists depicted not the visible but the actual sizes of angles $\alpha$,
$\beta$, $\alpha'$, $\beta'$, equal to 90°, producing depiction~C.  Perhaps in this
case, which should also be classed as reverse perspective, the distortion relative to
the visible form had the aim of conveying the actual rectangular silhouette not only
of the front and back board but of the book as a whole.

The ``primacy of continuity,'' which compelled the artist to give angles $\gamma$ and
$\delta$ not the visible but some compromise size, led also to a compromise depiction
(relative to the visible magnitudes of angles) of rectangular objects directed toward
the viewer by their narrow face.  Ill.\,43D gives a schematic depiction of a chest
arranged in this way.  The angles $\alpha$ and $\alpha'$ cannot be shown near to
right angles (as in scheme~B), since in that case the edge distant from the viewer
would become disproportionately large compared with the front edge (indeed their
visible sizes are very close, to a first approximation simply equal).  The only way
that allows one to show simultaneously the exact visible sizes of both angles and
edges is a break in the image; however, artists clearly considered the continuity of
the depiction of objects a very important property of the picture.  In the example
under consideration, taking the ``primacy of continuity'' into account, the most
reasonable representation appears to be some compromise depiction distributing the
distortions uniformly among all its components.

\srcpage{112}

A fragment of the icon ``Cosmas, Damian and James the Brother of God'' (Ill.\,44)
makes the foregoing arguments visually clear.  Here one sees how tactfully the artist
used the possibility afforded him by the system of perceptual perspective to
distribute unavoidable distortions reasonably.  The comparatively shallow depiction
of the Gospel book the artist considered it necessary to give with locally correct
transmission of the visible sizes of angles determining the silhouette, since the
small depth will not thereby lead to a disproportionately large increase of the edge
of the distant board relative to the near one.  The deeper depiction of the chest
compels giving the angles less near to right angles.  As a result the angle of
convergence of the edges~--- determining the degree of reverse perspective~--- turns
out to be approximately $-30°$ for the Gospel book and only approximately $5°$ for
the chest, while the relative increase of the distant edges compared with the near
ones in both cases is of the same order, approximately $10\%$.

\figblock{44}{``Cosmas, Damian and James the Brother of God.''  Fragment of an icon
  of the first half of the 16th century.  Museum of Architecture (MiAR).}

Thus, the form-constancy mechanism is capable under certain conditions of leading to
effects of reverse perspective.  These conditions are primarily: seeing objects
``from outside'' (when looking from the top and front) and the absence of breaks in
visual perception.  Acting only in relation to objects of familiar form, the
form-constancy mechanism weakens as the observed object recedes from the viewer.  The
effects described above can manifest themselves only in the depiction of space near
the artist.  The arguments presented here show that it is precisely medieval painting,
in particular Old Russian painting, with its heightened interest in the depiction of
near and shallow spaces, that was bound to give (and did give) numerous examples of
the types of reverse-perspective formation described here.

\figblock{43}{Schemes for depicting a Gospel book and a chest.}

\srcpage{113}

The discussion of the influence of the form-constancy mechanism on the peculiarities
of the depiction has shown that the attempt to convey one's visual perception runs up
against the impossibility of depicting volumetric bodies while preserving all the
geometrical features of the spatial image that has arisen in human consciousness.
Depending on the choice of what is ``primary'' and what is ``secondary,'' and for the
usual direction of the gaze, the depictions acquire an axonometric character or
acquire the geometrical forms of reverse perspective.  Whereas the size-constancy
mechanisms led to the result that under certain conditions a person simply \emph{saw}
in reverse perspective, the form-constancy mechanism produced reverse perspective in
the attempt to \emph{depict} the visible.

The natural question arises: can the form-constancy mechanism also sometimes lead to
actual vision in reverse perspective?  The answer must be affirmative.  For if the
person sees angles $\alpha$ and $\alpha'$ in Ill.\,43 as larger than in the retinal
image, and on the other hand visual perception is always characterised by the
continuity of objects, then all the reasoning carried through above retains its
force, and therefore in our vision too one may expect the perception of such objects
in reverse perspective.  From the author's experiments this is indeed the case when
observing from close range comparatively small parallelepipeds that fit entirely
within the field of sharp vision and are examined from one vertex.  True, these
features of visual perception are weakly expressed and are registered with greater
difficulty than the angles of ``divergence'' of the lines forming the linear
structure of the floor.  In this connection one may assert that although the
form-constancy mechanism does contribute to actual vision in reverse perspective, its
influence on the character of the depiction is sharply enhanced when one wishes to
convey on the picture plane one's perception of volumetric bodies without any
``distortions''~--- which, as was said above, is impossible.

\srcpage{114}

All four conditions for the appearance of reverse perspective examined here~--- in
the attempt to convey visual perception of near space on a picture~--- are of
entirely natural origin.  One may even assert that they correspond to normal human
visual perception.  At this point almost every reader must feel a quite legitimate
doubt~--- since ``he does not see it that way.''

Let us recall that the brain, which performs the work of transforming the retinal
image into the visible one, is highly susceptible to various kinds of suggestion.  In
Chapter~II it was already noted that the observer's set, communicated to him by
training, suggestion, or habit, can have a most powerful effect on perception.  The
modern person who, from children's books through photographs and cinema, sees almost
only images built according to the rules of linear perspective, has grown accustomed
to it and unconsciously tends always to see parallel lines converging as they recede
from the observer, which undoubtedly affects his perception.  To see parallel lines
diverging even slightly as they recede from the observer is possible only by
overcoming the ``training'' to which everyone has been subjected from
childhood.\footnote{It is, in general, sufficiently complex how things stand with
  the natural course of perception.  Much indicates that our perception of space is
  influenced by the perception of pictures.  The modern person, accustomed by
  photography and the pictures of new painting to direct perspective, ``fails to
  notice'' that much of what he sees is in slight reverse perspective.  It is quite
  reasonably conjectured that the regularities of vision are functions of the epoch,
  and it is not excluded that in time people may return to the natural perception of
  near objects in slight reverse perspective in everyday life.

  In the above-mentioned monograph \emph{Composition in Painting}, N.\,N.~Volkov
  asserts (p.~244) that ``divergence of parallels and increase of objects as they
  recede cannot be seen in the natural course of perception under any set.''  This
  erroneous assertion of N.\,N.~Volkov is probably connected with the fact that he
  conducted experiments on visual perception using the ``classical'' scheme of such
  experiments, which are never conducted by observing horizontal structures in
  foreshortening.}
That vision in slight reverse perspective is natural follows from the reaction of
people who were shown this on the example of, say, the linear structure of the floor:
no one spoke of unnaturalness of perception; all were surprised that they had never
noticed it before.  No less significant is the fact that afterwards people begin
constantly to see in reverse perspective, of course when the conditions necessary for
this are present.

\srcpage{115}

These ``necessary conditions'' were required by the medieval person as well.  As
already noted, they reduce to the small distance of the observed object from the
viewer, its relative smallness (corresponding to the region of sharp vision), and the
desirability of the appropriate foreshortening.  This last property of the geometry
of perceptual space leads to the result that, moving through a room, the person sees
the same object (for example, a footstool) now in axonometry, now in more or less
pronounced reverse perspective.  The geometric appearance of the object had to vary
depending on the foreshortening.  Reproducing his visual impressions subsequently,
the artist might remember that footstools are seen now in axonometry, now in slight
reverse perspective, and choose the type of depiction on the basis, for example, of
compositional considerations.  The medieval artist could see and depict objects in
slight reverse perspective, axonometry, and even in direct perspective (distant
planes).

What has been said here about the medieval artist's known indifference to strict
observance of some perspective system~--- in the ordinary sense of these words~--- can
be illustrated by many examples.  Thus, on the Pskov icons of the 14th century
``Cathedral of the Mother of God'' and ``Deesis,'' belonging to the brush of one
master,\footnote{L.\,Ya.~Ovchinnikov, N.\,V.~Kishilov.  Op.\,cit., ills.~15, 23.
  On the attribution of both icons to the brush of the ``Varvarinsky master'' see
  Catalogue, nos.~10, 12.}
the footstool is given in the first case in direct, in the second in reverse,
perspective.  It is notable that direct perspective was used by the artist for the
depiction of the footstool farther from the viewer, and reverse for the nearer one.

Moreover, one can find many examples where on one icon identical objects are depicted
in different ways.  On the ``Trinity'' of Andrei Rublev (Ill.\,33) the footstool of
the left angel (from the viewer's point of view) is given in slight reverse
perspective, while the footstool of the right one is in axonometry.  This peculiarity
of the geometry of the ``Trinity'' may stem from the already-mentioned
``indifference'' of the artist; but for such a reflective master as Rublev, an
observant man who strove to be precise in conveying the geometry of visual
perception, another explanation seems more probable.  Rublev may have noticed that
if a room contains, say, two footstools seen in different foreshortening, and if one
is perceived in axonometry, then at that moment the other appears in slight reverse
perspective (such an effect when observing objectively parallel lines has been
described above).  Perhaps it was precisely for this reason that, having depicted the
right angel's footstool in axonometry, he considered it necessary to give the left
footstool the geometry of reverse perspective.  Rublev's striving not to deviate
where possible from the geometry of visual perception follows from the analysis of
icons whose attribution to his brush is beyond doubt.  Especially carefully did the
master watch that the perspective constructions nowhere violate naturalness in the
rendering of objectively parallel lines.  On his icons perspective constructions
never give reverse perspective

\srcpage{116}

expressed more strongly than approximately 10°.  The footstool of the left angel of
the ``Trinity'' is characterised by a divergence angle of the edges of this order.
The upper edge of the Gospel book on the icon of the Apostle Paul from the iconostasis
of the Trinity Cathedral is depicted as near-axonometric or in reverse perspective
no stronger than 10° (this imprecision is connected with the fact that different
reproductions of the icon in albums give grounds for different assessments of this
angle); in any case this angle lies within natural limits.  It should be noted that
in connection with what has been said, the outer angle of the Gospel book (angle
$\alpha$ in Ill.\,43) is depicted probably more acute than the form-constancy
mechanism would call for.  Evidently Rublev considered that if one made this angle
as the form-constancy mechanism requires, a strong ``unnatural'' reverse perspective
would arise, which he, unlike many icon painters, considered impermissible.

Similarly the depiction of the Gospel book on the icon ``Apostle Paul'' from the
Vasily Tier.  Here two edges of the Gospel book are visible; the degree of
``reverseness'' of this depiction is somewhat stronger than in the analogous icon
from Zagorsk; assessing it by the two outer angles (angles $\alpha$ and $\alpha'$ in
Ill.\,43), the expression of reverse perspective is of the order of 25--30°.  It was
probably in order to weaken this effect that Rublev placed the apostle's hand on the
corner of the book in such a way as to ``break'' the continuity of the image of the
vivid vermilion edges and thereby soften the ``unnaturalness'' of the image.  This
concealment of the boundary between the two edges also simplified the depiction of
the clasps, which are shown as lines almost parallel to the upper boundary of the
edge.\footnote{For reproductions of these icons see, for example: M.~Alpátov.
  \emph{Andrei Rublev.}  Moscow, 1972, ills.~58, 70, 79.}
These observations speak of Rublev's striving not to deviate from naturalness in
perspective constructions.  This, it would seem, is contradicted by the form of the
throne in the ``Trinity'' (which Rublev, incidentally, left indeterminate, concealing
the edges of the throne with the angels' knees) and certain other peculiarities of
icons painted by him or by icon painters of his circle; but such a conclusion would
be erroneous, as will be shown later.

If we now return to the discussion of the problem of the appearance of reverse
perspective in art, with which the present chapter began, logic leads to the
conclusion that the problem does not really exist!  For this problem arose only
because the authors who worked on it were convinced of the impossibility of a person's
directly and vividly perceiving external space according to the laws of reverse
perspective.  Sometimes ingenious, sometimes overly abstruse theories, which were
meant to interpret and explain the appearance of reverse perspective in the visual
arts, were based on this conviction.\footnote{To show how firm this conviction was,
  let us cite the following: ``\ldots{}reverse perspective has nothing to do with
  the laws of genuine scientific perspective and encompasses early, conventional and
  imperfect methods of conveying real space on a plane.  In the practice of
  formalists, reverse perspective is already a deliberate means of distorting
  reality'' (\emph{Kratkiy slovar' terminov izobrazitel'nogo iskusstva} [Concise
  Dictionary of Terms of the Visual Arts], 4th ed.  Moscow, 1965, p.~121).}

\srcpage{117}

But this seemingly unassailable and obvious assertion~--- seemingly confirmed, moreover,
by the mathematical theory of linear perspective~--- turned out to be erroneous.

The fact of the appearance of reverse-perspective images of individual objects has an
absolutely simple explanation: people see that way.\footnote{It was repeatedly
  emphasised above that the various transformations of the retinal image characteristic
  of the visual perception process arose from biological needs and have as their aim
  an approximation of visual perception to the geometry of objective space.  From this
  standpoint, vision in reverse perspective would seem to have no justification.  The
  explanation of this apparent inconsistency is elementary: visual perception renders
  the main thing ``correctly'' (axonometrically), while slightly ``distorting'' the
  secondary (objects in foreshortening) in the direction of reverse perspective.  To
  avoid the latter, the unconscious work of the brain would have to be substantially
  complicated.  However, increasing the volume of the brain for the sake of such a
  secondary task is biologically unreasonable.  If the activity of the brain is not
  complicated, natural reverse perspective arises as an inevitable consequence of the
  axonometric character of basic spatial representation, the more tolerable in that it
  is small.}
Therefore one should be surprised not by the appearance of reverse perspective in
painting, but by its almost complete disappearance in modern art.

In conclusion of the present chapter let us return to the position stated at the end
of Chapter~III, regarding the possibility of a single system of scientific
perspective.  That position may now be reasonably refined: such a unified system of
scientific perceptual perspective takes on the character of linear perspective for
distant regions of space, and for near regions passes over into axonometry or slight
reverse perspective.
